\documentclass[11pt]{article}
\usepackage[margin=1in]{geometry}
\usepackage{amsmath,amssymb,microtype}
\usepackage[colorlinks=true,urlcolor=blue]{hyperref}
\title{Request for independent mathematical criticism}
\author{Jevint Felixciano\\Independent researcher\\
\small jevintfelixciano@protonmail.com}
\date{7 November 2026}
\begin{document}
\maketitle
\noindent\textbf{Manuscript:} Finite Menus of Spectral Subspaces:
Realizability and the Cost of Convexification. The v7 preparation changes
authorship, declarations, licensing, and packaging; the structural proofs
are unchanged from v6. No independent review or endorsement is claimed.

\paragraph{Requested assessment.}
Please challenge correctness and prior-art distinctness, rather than provide
a general endorsement. A counterexample, a missing hypothesis, or a precise
older theorem implying the result would be especially useful. The focused
reading is Theorem 4, its spectral face formula, Theorem 6, and Appendix D.
Numerical tests and the internal audit are not offered as proofs.

\paragraph{Model and notation.}
For real rank-$r$ orthogonal projectors $P$, shared symmetric reward $A$,
site costs $C_L$, overheads $d_L$, and $p\in[0,B]$, put
\[
 \alpha(P)=-\operatorname{tr}(AP),\quad
 c(P)=\min_L\{\operatorname{tr}(C_LP)+d_L\},\quad
 J(P,p)=\alpha(P)+pc(P),\quad F(p)=\min_PJ(P,p).
\]
The resource interpretation assumes positive-semidefinite matrices,
nonnegative overheads, and nonnegative scalar price. Write
$R(P,p)=J(P,p)-F(p)$ and
$e(a,b)=\min_P\max\{R(P,a),R(P,b)\}$.
Let $\mathcal L=\{(\alpha(P),c(P))\}$ and let $\mathcal G$ contain
\emph{all} projectors globally optimal at at least one price, including ties.
The $K$-menu regrets $E_K$, $\overline E_K$, and $E_K^{\mathcal G}$ use,
respectively, actual projectors, coefficient points in
$\operatorname{conv}\mathcal L$, and projectors in $\mathcal G$.

\paragraph{Theorem 4: exactness on every subinterval.}
For $a<b$ let $s=(F(b)-F(a))/(b-a)$ and
\[
 D(a,b)=\max_{p\in[a,b]}\{F(p)-F(a)-(p-a)s\},\qquad
 S(p)=\{c(P):J(P,p)=F(p)\}.
\]
The claim is that $e(a,b)=D(a,b)$ for \emph{every} subinterval if and
only if $S(p)=[\min S(p),\max S(p)]$ at every interior price.
This implies equality of the three menu values for every budget, but the
manuscript does not assert the converse from one fixed budget or domain.

Please examine: (i) the compactness and endpoint cases in the relaxed
oracle; (ii) constructing a chord of exactly a missing slope across a
corner; (iii) using attainment to obtain a strict, rather than merely
unattained, gap; and (iv) persistent ties and zero-regret cases.
For spectral data the claimed slope set is a union of closed cost intervals
from active cutoff eigenspaces. Please challenge real-projector recovery,
the active-site argument at $p>0$, and the single-site corollary.

\paragraph{Theorem 6: additional-preset repair.}
If $q$ interior prices have missing optimal slopes, the claim is
\[
 E^{\mathcal G}_{K+\min(K,q)}\leq\overline E_K\leq E_K.
\]
Without a finite $q$, the $2K$ replacement bound remains. Please examine
the interval-cover partition, boundary/singleton segments, and whether
charging a two-preset replacement to a strictly interior exceptional price
can ever charge the same price twice. This is deterministic selection of
a preset after observing price, not randomization.

Finiteness for each fixed spectral instance uses analytic eigenvalue
branches on an open neighborhood of the compact price domain, finitely
many rank-$r$ branch sums, removal of identical branches, and isolated
zeros of nonidentical analytic differences. Please check each hypothesis.
No efficient enumeration, polynomial bound, uniformly small $q$, or
sharpness of the additive term for all $K,q$ is claimed.

\paragraph{Closest-prior-art challenge.}
Appendix D explicitly acknowledges convex Pareto representatives, weighted
scalarization, tangent replacement, and finite-policy bounds. In particular,
convexifying an already selected menu does not change its linear value;
allowing its selected representatives to come from a larger convex hull can.
Please identify any earlier result that directly yields the all-subinterval
criterion or the exceptional-price refinement, with an explicit mapping
of assumptions and conclusions. Conversely, is the remaining refinement
mathematically useful enough to merit a paper? The companion prior-art map
provides the exact comparisons, not a claim of exhaustive priority.

\paragraph{Evidence and response.}
The artifact includes reproducible numerical checks and separate historical
validation reports. Their role is computational reproducibility, not
proof or external peer review. Please distinguish what you read, derived,
or executed. A short response organized as ``correctness'', ``prior art'',
``significance'', and ``required changes'' is sufficient; unreviewed parts
should be marked as such. There is no request for coauthorship, arXiv
endorsement, or a commitment to formal journal refereeing.

\paragraph{AI disclosure.}
AI assistance was used substantially in planning, literature search,
mathematical drafting, implementation, analysis, and manuscript preparation.
The author seeks independent criticism of precisely those central claims;
internal automated checks should not be treated as independent review.
\end{document}
